\section{Discussion}

This study frames single-cell perturbation prediction as a reliability modeling problem in public functional genomics data. A random split remains useful as a low-stress anchor, but it does not answer the same biological question as predicting unseen perturbations, low-support signatures, new datasets, or an independent public screen. The observed rank reversals show that a random-split winner should not be assumed to be the safest model for perturbation prioritization or pathway interpretation.

PTL addresses a different task from the perturbation predictor itself. A predictor estimates a response signature; PTL estimates whether that completed response is transportable enough to retain. This separation is useful because reliability decisions can be layered over transparent baselines and published model adapters without changing the expression predictor. In the present public-data benchmark, PTL substantially reduced false transportability relative to naive confidence and simpler deployment-feature baselines, indicating that support, novelty, split context, and model-output summaries contain actionable reliability information.

The failure-mode atlas adds biological interpretability to the reliability model. Severe failures were concentrated at novelty, low-support, and dataset-transfer boundaries, and case examples showed that high-confidence predictions can still reverse or degrade top-gene response structure. These patterns argue for reporting failure composition together with aggregate fidelity: a prediction system that improves average cosine in familiar settings may still be risky in exactly the regimes where computational inference is most attractive.

For practical reporting, perturbation-response studies should not rely on a single random-split score. We recommend reporting stress-split performance, false-transportability rate, risk--coverage behavior, and failure composition whenever predictions are used for pathway interpretation, perturbation retrieval, or prioritization. This makes reliability evidence visible before downstream biological claims are made.

PTL should be used when evaluating perturbation-response predictors under new perturbations, prioritizing predicted responses for pathway or gene-function interpretation, comparing models across split regimes, or deciding when to abstain from high-risk predictions. It is not a replacement for experimental confirmation. Rather, it is a public-data grounded reliability layer that helps decide which predicted signatures deserve downstream biological attention and which should be treated as transfer-boundary failures.
